%%%%%%%%%%%%%%%%%%%%%%%%%%%%%%%%%%%%%%%%%%%%%%%%%%%%%%%%%%%%%%%%%%%%%%
% Lớp 10 | Phương pháp tọa độ trong mặt phẳng (𝑂𝑥𝑦) | Ba đường conic trong mặt phẳng tọa độ
% Dạng 0H9?5-5: Phương trình chính tắc của hypebol
% Số câu: 7  (sinh tự động bởi tach_cau.py; từ nay thầy sửa trực tiếp file này)
%%%%%%%%%%%%%%%%%%%%%%%%%%%%%%%%%%%%%%%%%%%%%%%%%%%%%%%%%%%%%%%%%%%%%%

%%%%%%%%%%%%%%%% TRẮC NGHIỆM 4 PHƯƠNG ÁN %%%%%%%%%%%%%%%%

% nguồn: MaHoaNH_TracNghiem_Lop10
\begin{ex}%[0H9N5-5]%[ma:0b8178b6]
	Trong mặt phẳng $Oxy$, phương trình nào dưới đây là phương trình chính tắc của Hypebol? 
	\choice 
		 {$\dfrac{x^2}{25}+\dfrac{y^2}{16}=1$}
		{$\dfrac{x^2}{25}-\dfrac{y}{16}=1$} 
		 {$y^2=16x$}
		{\True $\dfrac{x^2}{25}-\dfrac{y^2}{16}=1$} 
	\loigiai{ 
		Ta có phương trình chính tắc  của Hypebol là $\dfrac{x^2}{25}-\dfrac{y^2}{16}=1$.
	 }
\end{ex}


%%%%%%%%%%%%%%%% ĐÚNG/SAI %%%%%%%%%%%%%%%%

% nguồn: MaHoaNH_DungSai_Lop10
\begin{ex}%[0H9H5-5]%[ma:9d93a118]
	Cho hypebol $(H)$ có tiêu cự bằng $20$ và hiệu các khoảng cách từ mỗi điểm thuộc hypebol đó tới hai tiêu điểm  đó bằng $12$. Trong các mệnh đề dưới đây, mệnh đề nào đúng, mệnh đề nào sai?
	\choiceTF 
		{Các tiêu điểm là $F_1(-10;0)$; $F_2(10;0)$} 
		{\True Phương trình chính tắc của hypebol là  $\dfrac{x^2}{36}-\dfrac{y^2}{64}=1$}
		{Các  giao điểm của $(H)$ với trục tung là $B_1(-8;0)$; $B_2(8;0)$}
		{\True Các giao điểm của $(H)$ với trục hoành là  $A_1(-10;0)$; $A_2(10;0)$}
	\loigiai{Gọi phương trình  chính tắc của hypebol là $\dfrac{x^2}{a^2}-\dfrac{y^2}{b^2}=1$ với $a>0$; $b>0$.\\
		Ta có  tiêu cự $2c=2\sqrt{a^2+b^2}=20$, suy ra $c=10$ và $2a=12$ nên $a=6$.\\
		Suy ra $b^2=c^2-a^2=100-36=64$.  Vậy phương trình chính tắc của hypebol là $\dfrac{x^2}{36}-\dfrac{y^2}{64}=1$.
	 }
\end{ex}

% nguồn: MaHoaNH_DungSai_Lop10
\begin{ex}%[0H9H5-5]%[ma:1be59ac3]
	Cho hypebol $(H)$ đi qua điểm $A(5;0)$ và  có một tiêu điểm là $F_2(9;0)$. Trong các mệnh đề dưới đây, mệnh đề nào đúng, mệnh đề nào sai?
	\choiceTF 
		 {\True Tiêu điểm $F_1(-9;0)$}
		{\True Hiệu các khoảng cách từ mỗi điểm trên hypebol tới hai tiêu  điểm có giá trị tuyệt đối bằng $10$}
		{Tiêu cự của hypebol  là $6$}
		{Hypebol có phương trình chính tắc là $\dfrac{x^2}{25}-\dfrac{y^2}{16}=1$} 
	\loigiai{Gọi phương trình  chính tắc của hypebol là $\dfrac{x^2}{a^2}-\dfrac{y^2}{b^2}=1$ với $a>0$; $b>0$.\\
		Ta có $A(5;0)\in  (H)$ nên $a^2=25$.\\
		 Lại có $c=9$, suy ra $b^2=c^2-a^2=56$.\\
		Vậy phương trình chính tắc của hypebol  là $\dfrac{x^2}{25}-\dfrac{y^2}{56}=1$.\\
		Tiêu cự của hypebol là  $F_1F_2=2c=18$.
	 }
\end{ex}

% nguồn: MaHoaNH_DungSai_Lop10
\begin{ex}%[0H9H5-5]%[ma:9ead55fa]
	Cho hypebol $(H)$ có phương trình chính tắc là $\dfrac{x^2}{a^2} - \dfrac{y^2}{b^2} = 1$, biết rằng  $(H)$ đi qua điểm $M\left(\sqrt{7};0\right)$ và $N\left(2\sqrt{7};3\sqrt{3}\right)$. Trong các mệnh đề dưới đây, mệnh đề nào đúng, mệnh đề nào sai?
	\choiceTF 
		{\True Tiêu  cự là $8$}
		{\True  $\dfrac{28}{a^2}-\dfrac{27}{b^2}=1$}
		{\True Phương trình chính tắc  của hypebol là $\dfrac{x^2}{7}-\dfrac{y^2}{9}=1$}
		{Hiệu các khoảng cách từ một  điểm trên hypebol tới hai tiêu điểm có giá trị tuyệt đối bằng $4\sqrt{7}$}
	 \loigiai{
		Gọi phương trình chính tắc của hypebol là $\dfrac{x^2}{a^2}-\dfrac{y^2}{b^2}=1$  với $a>0$; $b>0$.\\
		Ta  có hệ phương trình
		$$\heva{&\dfrac{7}{a^2}=1\\&\dfrac{28}{a^2}-\dfrac{27}{b^2}=1}\Leftrightarrow  \heva{&a^2=7\\&b^2=9.}$$
		Phương trình chính  tắc của hypebol là $\dfrac{x^2}{7}-\dfrac{y^2}{9}=1$.\\
		Độ  dài tiêu cự là $2c=2\sqrt{a^2+b^2}=8$.\\
		Hiệu các khoảng  cách từ một điểm trên hypebol tới hai tiêu điểm có giá trị tuyệt đối bằng $2a=2\sqrt{7}$.
	} 
\end{ex}

% nguồn: MaHoaNH_DungSai_Lop10
\begin{ex}%[0H9H5-5]%[ma:10abd227]
	Cho hypebol $(H)$ có phương trình chính tắc là $\dfrac{x^2}{a^2} - \dfrac{y^2}{b^2} = 1$, biết rằng  $(H)$ đi qua điểm $M\left(3\sqrt{2};4\right)$ và hiệu các khoảng cách từ một điểm trên hypebol tới hai tiêu điểm có giá trị tuyệt đối bằng $6$. Trong các mệnh đề dưới đây, mệnh đề nào đúng, mệnh đề nào sai?
	 \choiceTF
		{Độ  dài tiêu cự là $2c=2\sqrt{a^2+b^2}=12$}
		 {\True $\dfrac{18}{a^2}-\dfrac{16}{b^2}=1$}
		{\True  Phương trình chính tắc của hypelbol là $\dfrac{x^2}{9}-\dfrac{y^2}{16}=1$}
		{Hiệu các khoảng cách từ một điểm trên hypebol tới  hai tiêu điểm có giá trị tuyệt đối bằng $8$}
	 \loigiai{
		Gọi phương trình chính tắc của  hypelbol là $\dfrac{x^2}{a^2}-\dfrac{y^2}{b^2}=1$ với $a>0$; $b>0$.\\
		Ta có hệ phương  trình
		$$\heva{&\dfrac{18}{a^2}-\dfrac{16}{b^2}=1\\&2a=6}\Leftrightarrow \heva{&a^2=9\\&b^2=16.}$$ 
		Phương trình chính tắc của hypebol  là $\dfrac{x^2}{9}-\dfrac{y^2}{16}=1$.\\
		Độ dài tiêu cự là $2c=2\sqrt{a^2+b^2}=10$.\\ 
		Hiệu các khoảng cách từ một điểm trên hypebol tới hai tiêu điểm có giá trị tuyệt đối  bằng $2a=6$.
	 }
\end{ex}


%%%%%%%%%%%%%%%% TỰ LUẬN %%%%%%%%%%%%%%%%

% nguồn: MaHoaNH_TuLuan_Lop10
\begin{bt}%[0H9H5-5]%[ma:67dcba11]
	Viết phương trình chính tắc của một hyperbol  có các tiêu điểm $F_1(-5;0)$, $F_2(5;0)$ và đi qua điểm $M(6;2)$.
	 \loigiai{
		Hai tiêu điểm nằm trên trục hoành và đối xứng qua gốc tọa độ, nên hyperbol  có
		\begin{itemize}
			\item Tâm $O(0;0)$. 
			\item Trục thực nằm trên trục  $Ox$.
			\item Dạng  phương trình $\dfrac{x^2}{a^2} - \dfrac{y^2}{b^2} = 1$.
		\end{itemize}
		 Khoảng cách giữa hai tiêu điểm $2c = 10 \Rightarrow c = 5 \Rightarrow c^2 = 25$.\\
		Áp dụng công  thức $c^2 = a^2 + b^2 \Rightarrow b^2 = c^2 - a^2 = 25 - a^2$.\\
		 Thay điểm $M(6;2)$ vào phương trình
		 \[
		\dfrac{6^2}{a^2} -  \dfrac{2^2}{25 - a^2} = 1
		\Rightarrow \dfrac{36}{a^2} -  \dfrac{4}{25 - a^2} = 1.
		\] 
		Giải phương  trình
		\[ 
		 \dfrac{36(25 - a^2) - 4a^2}{a^2(25 - a^2)} = 1
		\Rightarrow \dfrac{900 -  36a^2 - 4a^2}{a^2(25 - a^2)} = 1
		\Rightarrow \dfrac{900 - 40a^2}{a^2(25 - a^2)} =  1
		\Rightarrow a^4 - 65a^2  + 900 = 0.
		\] 
		Đặt $t =  a^2$, ta có
		\[ 
		t^2 - 65t + 900 = 0  \Rightarrow t = 45 \text{ hoặc } t = 20.
		 \]
		Xét 
		\begin{itemize}
			\item Với $a^2  = 45 \Rightarrow b^2 = 25 - 45 = -20$ (vô lý, loại).
			\item Với $a^2 = 20 \Rightarrow b^2 = 25 - 20 =  5 > 0$ (chọn).
		\end{itemize}
		Vậy phương  trình hyperbol cần tìm là $\dfrac{x^2}{20} - \dfrac{y^2}{5} = 1$.
	 }
\end{bt}

% nguồn: MaHoaNH_TuLuan_Lop10
\begin{bt}%[0H9V5-5]%[ma:9d9e6bb4]
	Viết phương trình chính tắc  của elip có tiêu điểm là $F_1(-4;0)$, $F_2(4;0)$ và đi qua điểm $M(3;2)$.
	 \loigiai{
		Vì hai tiêu điểm $F_1(-4;0)$, $F_2(4;0)$ nằm trên  trục hoành, nên elip có trục lớn nằm trên trục $Ox$ và tâm $O(0;0)$.\\
		Khoảng  cách giữa hai tiêu điểm
		\[ 
		2c  = \mathrm{d}(F_1, F_2) = 8 \Rightarrow c = 4
		\Rightarrow  c^2 = 16.
		 \]
		Phương  trình elip dạng
		\[ 
		 \dfrac{x^2}{a^2} + \dfrac{y^2}{b^2} = 1, \quad \text{với } c^2 = a^2 - b^2 \Rightarrow b^2 = a^2 - 16.
		 \]
		Vì điểm $M(3;2)$ thuộc elip nên  thỏa mãn phương trình
		 \[
		\dfrac{3^2}{a^2}  + \dfrac{2^2}{a^2 - 16} = 1
		\Rightarrow \dfrac{9}{a^2}  + \dfrac{4}{a^2 - 16} = 1.
		\] 
		Giải phương  trình
		\[ 
		 \dfrac{9(a^2 - 16) + 4a^2}{a^2(a^2 - 16)} = 1
		\Rightarrow \dfrac{9a^2  - 144 + 4a^2}{a^4 - 16a^2} = 1
		 \Rightarrow \dfrac{13a^2 - 144}{a^4 - 16a^2} = 1.
		 \]
		Suy  ra
		 \[
		13a^2 - 144 =  a^4 - 16a^2
		\Rightarrow a^4  - 29a^2 + 144 = 0.
		\] 
		Đặt $t  = a^2$, phương trình trở thành
		\[ 
		t^2 - 29t + 144  = 0 \Rightarrow t = 16 \text{ hoặc } t = 9.
		\] 
		Với $a^2 = 16 \Rightarrow  b^2 = 0$ (không phải elip), loại.\\
		Với $a^2 = 9 \Rightarrow b^2 = 9 - 16 = -7$ (vô lý), loại.  \\
		Vậy không tồn  tại elip có hai tiêu điểm $F_1(-4;0)$, $F_2(4;0)$ và đi qua $M(3;2)$.\\
		\textbf{Kết luận} Không tồn  tại elip thỏa mãn các điều kiện đã cho.
	} 
\end{bt}
